\documentclass[11pt]{article}
\usepackage[margin=1in]{geometry}
\usepackage{amsmath,amssymb,amsthm}
\usepackage{hyperref}

\newtheorem{lemma}{Lemma}
\newtheorem{proposition}{Proposition}
\theoremstyle{definition}
\newtheorem{remark}{Remark}

\title{A Disproof of Conjecture \texttt{00000001231}\\
\large On the Jacobian order of the Johnson graph $J(n,k)$}
\author{TLMC Submission}
\date{September 2026}

\begin{document}
\maketitle

\begin{abstract}
Conjecture \texttt{00000001231} asserts that the order of the Jacobian (spanning-tree
torsion group) of the Johnson graph $J(n,k)$ equals
\[
|\operatorname{Jac}(J(n,k))| \;=\; \binom{n-2}{k-1}^{\binom{n}{k}-1}\cdot \prod_i f(i),
\qquad f(i) = (i^2-i+1)^{e_i},
\]
for some exponents $e_i$. We show that the stated formula is \emph{false as written}:
every factor $i^2-i+1$ is odd, so the product term is always odd, whereas at the
admissible point $(n,k)=(4,1)$ the true value of the Jacobian order is
$\tau(K_4)=16$, which is even. The matrix-tree theorem itself is used as the measuring
device and is not disputed.
\end{abstract}

\section{Setup}

For a finite connected graph $G$, the \emph{Jacobian} (a.k.a.\ sandpile group) of $G$
is a finite abelian group whose order equals the number $\tau(G)$ of spanning trees of
$G$. By the \emph{matrix-tree theorem},
\begin{equation}\label{eq:mtt}
\tau(G) \;=\; \det L^{\circ},
\end{equation}
where $L$ is the graph Laplacian and $L^{\circ}$ is any $(|V|-1)\times(|V|-1)$ principal
minor obtained by deleting one row and the corresponding column.

The \emph{Johnson graph} $J(n,k)$ has as vertices the $k$-subsets of $\{1,\dots,n\}$,
with two vertices adjacent exactly when their intersection has size $k-1$ (equivalently,
their symmetric difference has size $2$).

\begin{lemma}\label{lem:Jn1}
$J(n,1)$ is the complete graph $K_n$.
\end{lemma}
\begin{proof}
Every $1$-subset $\{x\}$ intersects every other $1$-subset $\{y\}$, $x \neq y$, in the
empty set, which has size $1-1=0=k-1$; hence all pairs of distinct vertices are
adjacent, and there are $\binom{n}{2}$ edges. This is exactly $K_n$.
\end{proof}

\section{The counterexample $(n,k)=(4,1)$}

By Lemma~\ref{lem:Jn1}, $J(4,1) = K_4$. Its Laplacian and the relevant $3\times 3$
principal minor are
\[
L(K_4) \;=\;
\begin{pmatrix}
 3 & -1 & -1 & -1\\
-1 &  3 & -1 & -1\\
-1 & -1 &  3 & -1\\
-1 & -1 & -1 &  3
\end{pmatrix},
\qquad
L^{\circ} \;=\;
\begin{pmatrix}
 3 & -1 & -1\\
-1 &  3 & -1\\
-1 & -1 &  3
\end{pmatrix}.
\]
Expanding $\det L^{\circ}$ along the first row:
\begin{align*}
\det L^{\circ}
 &= 3\cdot\det\begin{pmatrix}3&-1\\-1&3\end{pmatrix}
   -(-1)\cdot\det\begin{pmatrix}-1&-1\\-1&3\end{pmatrix}
   +(-1)\cdot\det\begin{pmatrix}-1&3\\-1&-1\end{pmatrix}\\[2pt]
 &= 3\,(9-1) \;+\; 1\,(-3-1) \;-\; 1\,(1+3)\\
 &= 24 \;-\; 4 \;-\; 4\\
 &= 16.
\end{align*}
This agrees with the Cayley--Pr\"ufer value $\tau(K_n)=n^{n-2}=4^2=16$. Hence
\begin{equation}\label{eq:tauK4}
|\operatorname{Jac}(J(4,1))| \;=\; \tau(K_4) \;=\; 16 .
\end{equation}

\section{The parity obstruction}

\begin{lemma}[Parity of the small factors]\label{lem:odd}
For every integer $i$, the number $i^2 - i + 1$ is odd.
\end{lemma}
\begin{proof}
$i^2 - i = i(i-1)$ is a product of two consecutive integers, one of which is even;
hence $i^2-i$ is even and $i^2-i+1$ is odd.
\end{proof}

\begin{proposition}[The conjecture fails at $(4,1)$]\label{prop:main}
There are no nonnegative integers $e_i$ such that
\[
\binom{2}{0}^{\binom{4}{1}-1}\cdot\prod_i (i^2-i+1)^{e_i} \;=\; 16 .
\]
\end{proposition}
\begin{proof}
The leading factor is $\binom{2}{0}^{3}=1$, so the conjecture asserts
$\prod_i (i^2-i+1)^{e_i} = 16$. By Lemma~\ref{lem:odd} each factor $i^2-i+1$ is odd,
and any product of odd integers is odd. But $16$ is even. Contradiction.
\end{proof}

\section{A second instance}

$J(4,2)$ is the line graph $L(K_4)$, i.e.\ the octahedral graph, with $6$ vertices and
$12$ edges. Deleting the first row and column of its Laplacian and evaluating the
resulting $5\times5$ determinant (carried out exactly in the accompanying
\texttt{reproduce.py}) gives
$\tau(J(4,2)) = 384 = 2^7 \cdot 3$. The conjecture's leading factor at $(4,2)$ is
$\binom{2}{1}^{\binom{4}{2}-1} = 2^5 = 32$, leaving a quotient $384/32 = 12$, which is
again even --- impossible for a product of odd factors $i^2-i+1$. The parity obstruction
is thus not an artifact of a single point.

\section{Scope of the disproof}

\begin{remark}
We refute only the \emph{literal} formula of Conjecture \texttt{00000001231}, at the
concrete admissible point $(n,k)=(4,1)$ (with corroboration at $(4,2)$). The
matrix-tree theorem, invoked by the conjecture itself, is of course correct and is used
here only as a tool. We make no claim about what corrected formula, if any, the author
intended.
\end{remark}

\section{Machine verification}

The directory \texttt{lean4/} contains a Lean~4 project (core only, toolchain
\texttt{leanprover/lean4:v4.33.1}) that formally establishes, with zero axioms and no
\texttt{sorry}:
\begin{itemize}
\item \texttt{Tlmc1231.tau\_K4}: the $3\times3$ integer determinant above equals $16$;
\item \texttt{Tlmc1231.f\_odd}: $\forall\, i : \mathbb{N}$, $(i\cdot i - i + 1) \bmod 2 = 1$;
\item \texttt{Tlmc1231.prod\_f4}: $\prod_{i<4}(i^2-i+1) = 21$, together with
      \texttt{prod\_f4\_odd} and \texttt{sixteen\_even};
\item \texttt{Tlmc1231.refute}: the parity-incompatible instance
      $\neg\bigl(\prod_{i<4}(i^2-i+1) = 16 \wedge \tau = 16\bigr)$.
\end{itemize}
The script \texttt{reproduce.py} recomputes $\tau(K_4)=16$, $\tau(J(4,2))=384$, the
leading factors $1$ and $32$, and the parity of $i^2-i+1$ for $0 \le i \le 100$.

\end{document}
